\section{解凍アルゴリズムの説明}
最初にoutputDirが与えられたかどうかを判断する。与えられた場合、そのoutputDirが実在するかどうかを判断する。全部満たす場合のみDestinationをprint outする。その後、archiveファイルを開き、そこに保存された圧縮後のファイルを一つずつdecompressFileFromArchive(archive, outputDir, 1)を用いて解凍する。decompressFileFromArchive関数では、archiveファイル先頭にあるfrequencyListを読み取り、圧縮時と同様にHuffman Treeを再構築する。手法が同じなため構築されたHuffman Treeも同じである。左を0、右を1とし、圧縮されたファイルから0もしくは1を一つずつ読み込み、rootから辿ると最終的に到達したnodeが対応するcharをinputFileに書き込む。それを繰り返すことで最終的には解凍されたファイルのサイズが記録されたFileHeaderのoriginalSizeと等しくなり、つまり圧縮されたファイルの一つが完全に解凍される。archive fileの先頭にファイルの総数も書かれてあるため、それに応じて循環の回数を決めれば最終的には全部のファイルが解凍される。

\end{verbatim}